
encore int\`egre sur @@M0@@, \`a fibres de dimension @@M1@@\,; en vertu de 4.3, on peut remplacer @@M2@@ par @@M3@@. Enfin, en vertu de 4.3 et 3.10 (utilisant la r\'esolution des singularit\'es des surfaces), quitte \`a remplacer @@M4@@ par un rev\^etement plat d'un ouvert non-vide de @@M5@@, on peut remplacer @@M6@@ par un @@M7@@-sch\'ema lisse.

Maintenant, gr\^ace \`a 4.6 (i) @@M8@@ (iii) et au th\'eor\`eme de finitude et de sp\'ecialisation des groupes de cohomologie \'etale (SGA 4 XVI 2.2), 5.1 r\'esulte de la proposition suivante.

\textbf{Proposition 5.2.} \emph{Soient @@M9@@ un corps alg\'ebriquement clos et @@M10@@ un @@M11@@-sch\'ema propre et lisse. Alors le groupe @@M12@@ des cycles alg\'ebriques de codimension @@M13@@ sur @@M14@@, modulo l'\'equivalence num\'erique, est libre de rang @@M15@@, le nombre de Betti.}

En effet, si l'on oublie l'action des racines de 1, la cohomologie \'etale @@M16@@ o\`u @@M17@@, devient une cohomologie de Weil~: elle satisfait \`a la dualit\'e de Poincar\'e et \`a la formule de K\"unneth\,; et elle re\c{c}oit un homomorphisme fonctoriel, non nul de l'anneau des cycles alg\'ebriques modulo l'\'equivalence rationnelle @@M18@@ (voir X 7.13).

Soient @@M19@@, o\`u @@M20@@, des \'el\'ements qui s'en vont sur une base du @@M21@@-espace vectoriel dans @@M22@@, engendr\'e par les cycles alg\'ebriques. Consid\'erons l'homomorphisme
@@M23@@
d\'efini par @@M24@@, o\`u @@M25@@ d\'esigne le nombre d'intersection. Alors on voit facilement que l'image de @@M26@@ est @@M27@@.

